% =========================================================
\section{Experimental Results}
\label{sec:results}
% =========================================================

\subsection{Main Retrieval Results}

Table~\ref{tab:retrieval_results} reports retrieval performance on the test split.
TEF-RAG achieves the highest value on all four primary metrics: Recall@5, nDCG@5, Complete@5, and FlowComplete@5.
Compared with the strongest non-TEF baseline on these metrics, which is Temporal-BM25 in all cases,
TEF-RAG improves Recall@5 from 0.4591 to 0.7228 and nDCG@5 from 0.4955 to 0.6231.
The difference is larger for set-level completeness: Complete@5 increases from 0.0792 to 0.3188, and FlowComplete@5 from 0.0792 to 0.3146.
These results indicate that adding temporal preference alone can improve retrieval of individual evidence records, but explicitly modeling evidence-flow relations and jointly selecting evidence sets is more important for recovering a complete task-support structure.

\begin{table}[t]
\caption{Main retrieval results on the test split (Top-5).}
\label{tab:retrieval_results}
\centering
\small
\setlength{\tabcolsep}{3pt}
\resizebox{\columnwidth}{!}{%
\begin{tabular}{lrrrr}
\toprule
Method & Recall & nDCG & Complete & FlowComplete \\
\midrule
BM25 & 0.3791 & 0.3915 & 0.0563 & 0.0563 \\
BGE Reranker & 0.2916 & 0.3070 & 0.0417 & 0.0417 \\
Temporal-BM25 & 0.4591 & 0.4955 & 0.0792 & 0.0792 \\
TA-RAG & 0.2769 & 0.2907 & 0.0375 & 0.0375 \\
\textbf{TEF-RAG} & \textbf{0.7228} & \textbf{0.6231} & \textbf{0.3188} & \textbf{0.3146} \\
\bottomrule
\end{tabular}%
}
\end{table}

\subsection{Validation Ablation Study}
\label{sec:ablation_results}

Table~\ref{tab:ablation_results} reports the four ablations on the validation split. Full TEF-RAG achieves a FlowComplete@5 of 0.4042. Removing learned pair proposal reduces it only to 0.4021, indicating that under a fixed pair budget, the proposal module contributes relatively little directly to final retrieval quality; its main value is to focus relation reasoning on more promising pairs. In contrast, removing relation-graph-derived features reduces FlowComplete@5 to 0.3667, showing that the set ranker does make use of graph-structural information.

The training objective and the ranker itself have a larger effect. Removing hard-negative mining reduces FlowComplete@5 from 0.4042 to 0.2833, while removing the nonlinear set ranker and falling back to the Stage-3A selector reduces it to 0.3125. These results indicate that final performance is not explained only by candidate recall or additional relation reasoning; it depends strongly on explicit learning from high-relevance but process-incomplete hard negatives and on nonlinear set-level ranking. Notably, removing the nonlinear ranker slightly increases nDCG@5 relative to the full model (0.6965 vs.\ 0.6885), while substantially decreasing Complete@5 and FlowComplete@5. This further supports that the method optimizes process completeness at the set level rather than conventional relevance ranking alone.

\begin{table}[t]
\caption{TEF-RAG ablation results on the validation split.}
\label{tab:ablation_results}
\centering
\small
\setlength{\tabcolsep}{3pt}
\resizebox{\columnwidth}{!}{%
\begin{tabular}{lrrrr}
\toprule
Variant & Recall & nDCG & Complete & FlowComplete \\
\midrule
Full TEF-RAG & 0.7567 & 0.6885 & 0.4063 & 0.4042 \\
w/o Pair Proposal & 0.7494 & 0.6832 & 0.4063 & 0.4021 \\
w/o Graph Features & 0.7184 & 0.6601 & 0.3750 & 0.3667 \\
w/o Hard Negatives & 0.6748 & 0.6042 & 0.2875 & 0.2833 \\
w/o Nonlinear Ranker & 0.7070 & 0.6965 & 0.3188 & 0.3125 \\
\bottomrule
\end{tabular}%
}
\end{table}

\subsection{Downstream Structured Generation Results}
\label{sec:generation_results}

Table~\ref{tab:generation_results} reports the four generation metrics that most directly answer whether better evidence flow improves the final work order and action plan.
Under the same generator, prompt, schema, and decoding settings, TEF-RAG achieves the highest value on all four metrics:
Field Macro F1 of 0.5345, Action F1 of 0.2035, Citation F1 of 0.1585, and Plan EM (strict) of 0.0417.
Compared with the strongest non-TEF baseline for each metric, TEF-RAG yields relative improvements of approximately 24.3\%, 26.9\%, and 121.8\% on Action F1, Citation F1, and Plan EM, respectively; Field Macro F1 improves by approximately 2.4\% over Temporal-BM25.
These results show that the benefit of TEF-RAG is not limited to upstream retrieval hit rates, but transfers to the quality of key work-order fields, maintenance actions, and evidence citations.
TEF-RAG improves field-level accuracy, action correctness, and plan-level consistency in the generated work orders.

\begin{table}[t]
\caption{Primary structured-generation metrics on the generation test set.}
\label{tab:generation_results}
\centering
\small
\setlength{\tabcolsep}{3pt}
\resizebox{\columnwidth}{!}{%
\begin{tabular}{lrrrr}
\toprule
Method & Field F1 & Action F1 & Citation F1 & Plan EM \\
\midrule
BM25 & 0.4715 & 0.1372 & 0.0813 & 0.0167 \\
BGE Reranker & 0.4129 & 0.1198 & 0.0405 & 0.0146 \\
Temporal-BM25 & 0.5218 & 0.1637 & 0.1249 & 0.0167 \\
TA-RAG & 0.4059 & 0.1115 & 0.0379 & 0.0188 \\
\textbf{TEF-RAG} & \textbf{0.5345} & \textbf{0.2035} & \textbf{0.1585} & \textbf{0.0417} \\
\bottomrule
\end{tabular}%
}
\end{table}

% =========================================================
\section{Limitations and Threats to Validity}
\label{sec:limitations}
% =========================================================

The benchmark in this work is a public-source-grounded, AI-assisted synthetic benchmark rather than a collection of real power-station work orders or domain-expert annotations.
Its primary value therefore lies in controlled comparison of temporal evidence-retrieval mechanisms; the results should not be directly extrapolated to real-world power-station fault frequencies or field robotic safety.
Although the generation evaluation covers fields, actions, dependencies, and citations, it does not simulate actual robotic execution, environmental feedback, or safety interlocks.

% =========================================================
\section{Conclusion}
\label{sec:conclusion}
% =========================================================

We presented TEF-RAG, which extends the retrieval objective for dynamic O\&M RAG from independent document relevance to a Temporal Evidence Flow that is visible at the query time, relation-consistent, and process-complete.
On the frozen temporal-hard benchmark, TEF-RAG outperforms four comparison methods on Recall@5, nDCG@5, Complete@5, and FlowComplete@5.
In the structured generation test with a fixed downstream generator, TEF-RAG also achieves the highest field-level match, Action F1, Citation F1, and strict Plan EM.
These results indicate that retrieval over evidence sets and their temporal relations can improve the task-support context supplied to the generator and, in turn, increase the accuracy of work-order fields, maintenance actions, and evidence citations.
Future work should further integrate explicit planning and execution constraints.
